\chapter{Bose-Einstein-Kondensat (BEC)-Objekte}

\section{Einleitung}

In der WDBT+ werden die kompaktesten Objekte des Universums nicht als Schwarze Löcher mit Singularitäten beschrieben, sondern als Bose-Einstein-Kondensate (BEC) aus Cooper-Paaren. Diese BEC-Objekte sind singularitätsfrei, ihre minimale Größe wird durch die fundamentale Längenskala \( l_0 \) begrenzt, und ihre maximale Masse folgt direkt aus der fraktalen Raumdimension \( D \).

Die beobachteten supermassereichen Objekte in Galaxienzentren (bis \( 10^{10} M_\odot \)) sind keine einzelnen kompakten Objekte, sondern Systeme aus einem zentralen BEC-Objekt und einer massiven Umgebung (Akkretionsscheibe, Sternhaufen, Gas). Diese Umgebung entsteht durch kontinuierliche Materieentstehung im Quantenpotential des Kerns und die anschließende Drift nach außen.

\section{Die fundamentale Längenskala \( l_0 \)}

Aus der fraktalen Raumstruktur mit Dimension \( D \) folgt eine fundamentale Längenskala:

\[
l_0 = \left( \frac{V_{\text{Dodekaeder}}}{V_{\text{Einheitskugel}}} \right)^{1/3} \lambda_p \approx 1,8 \times 10^{-15} \text{ m},
\]

wobei \( \lambda_p \) die Compton-Wellenlänge des Protons ist. Diese Länge ist die kleinste physikalisch mögliche Distanz – nichts kann kleiner sein. Sie ersetzt die Planck-Skala als natürliche untere Grenze.

\section{Paarbildung und Bose-Einstein-Kondensation}

Bei hinreichend hohen Dichten bilden Fermionen Cooper-Paare und werden zu effektiven Bosonen. Diese können in einen Bose-Einstein-Kondensat (BEC)-Zustand übergehen. In diesem Zustand:

\begin{itemize}
    \item Entfällt der Fermi-Druck (der sonst weitere Kompression verhindert)
    \item Das Quantenpotential dominiert die repulsive Wechselwirkung
    \item Das Objekt kann bis zur fundamentalen Länge \( l_0 \) schrumpfen
\end{itemize}

\section{Zustandsgleichung eines BEC-Objekts}

Aus der Gleichgewichtsbedingung zwischen Gravitation und Quantenpotential im fraktalen Raum folgt:

\[
P(\rho) = K \cdot \rho^\gamma
\]

mit

\[
\gamma = \frac{3}{2}, \quad K = \frac{\hbar^2}{2m_B} \left( \frac{3\pi^2}{g} \right)^{2/3} \cdot \frac{1}{m_B^{2/3}} \cdot f(D)
\]

wobei \( f(D) \approx 1,2 \) für \( D \approx 2,704 \) und \( m_B \approx 2m_p \) die Masse eines Cooper-Paares ist.

\section{Radius-Masse-Beziehung}

Aus der Gleichgewichtsbedingung \( dE/dR = 0 \) mit

\[
E(R) = -C_G \frac{GM^2}{R^{D-2}} + A \frac{\hbar^2 M}{m_B^2 R^{2D/3}}
\]

folgt durch Ableitung und Auflösen nach \( R \):

\[
R^{2-D/3} = \frac{2DA\hbar^2}{3(D-2)C_G G m_B^2} \cdot \frac{1}{M}
\]

Mit \( D \approx 2,704 \) ergibt sich \( 2 - D/3 \approx 1,0948 \), also:

\[
R(M) = K_R \cdot M^{-0,9134}
\]

Numerisch: \( K_R \approx 5,7 \times 10^{13} \, \text{m} \cdot \text{kg}^{0,9134} \).

\section{Maximale Masse eines BEC-Objekts}

Die minimale Größe ist \( R_{\text{min}} = l_0 \). Daraus folgt:

\[
M_{\text{BEC}} = \left( \frac{K_R}{l_0} \right)^{1/0,9134} \approx 3 \times 10^6 M_\odot
\]

Dies ist die maximale Masse eines einzelnen BEC-Objekts. Oberhalb dieser Grenze wäre \( R < l_0 \), was physikalisch unmöglich ist.

\section{Relativistische Korrekturen aus der Weber-Gravitation}

Die Weber-Gravitation führt zu einem modifizierten innersten stabilen Orbit (ISCO). Für eine kreisförmige Bahn (\( \dot{r} = 0, \ddot{r} = -v^2/r \)) liefert die Weber-Gravitation mit \( \beta = 0,5 \):

\[
\frac{mv^2}{r} = \frac{GMm}{r^2} \left( 1 + \frac{v^2}{2c^2} \right)
\]

Auflösen nach \( v^2 \):

\[
v^2 = \frac{GM}{r} \left( 1 + \frac{GM}{2c^2 r} \right)
\]

Die Bedingung für den innersten stabilen Orbit (\( dv/dr = 0 \)) ergibt:

\[
r_{\text{ISCO}} = \frac{12GM}{c^2}
\]

Für \( M = M_{\text{BEC}} \) ergibt sich \( r_{\text{ISCO}} \approx 5,3 \times 10^{10} \, \text{m} \approx 0,35 \, \text{AU} \) – makroskopisch.

\section{Beobachtete supermassereiche Objekte}

Beobachtet werden zentrale Objekte in Galaxien mit Massen bis \( 10^{10} M_\odot \). Diese sind keine einzelnen kompakten Objekte, sondern Systeme bestehend aus:

\begin{itemize}
    \item Einem zentralen BEC-Objekt mit \( M = M_{\text{BEC}} \approx 3 \times 10^6 M_\odot \) und \( R = l_0 \)
    \item Einer massiven Umgebung (Akkretionsscheibe, Sternhaufen, Gas), die die zusätzliche Masse von \( 10^9 - 10^{10} M_\odot \) enthält
\end{itemize}

Diese Umgebung entsteht durch die kontinuierliche Materieentstehung aus dem Quantenpotential des Kerns und die anschließende Drift nach außen.

\section{Strukturelle Analogie zum Atom}

Die beschriebene Struktur – ein winziger, extrem dichter Kern, umgeben von einer riesigen, dünnen Hülle, mit viel Vakuum dazwischen – ist analog zum Atom. Im Atom sitzt fast die gesamte Masse im Kern (\( \sim 10^{-15} \) m), die Elektronen sind weit draußen (\( \sim 10^{-10} \) m). In der BEC-Galaxie ist der Kern ebenfalls winzig (\( l_0 \approx 10^{-15} \) m), die Umgebung (Sterne, Gas) erstreckt sich über \( 10^{21} \) m – und der Kern trägt nur einen kleinen Teil der Gesamtmasse.

Die fraktale Raumdimension \( D \approx 2,704 \) verbindet beide Welten – sie bestimmt sowohl die Feinstrukturkonstante als auch das Massenverhältnis \( M_{\text{Gal}}/M_{\text{Kern}} \approx 1000 \).

\section{Kollisionen von BEC-Objekten}

BEC-Objekte haben keine Sonderstellung. Bei Kollisionen verhalten sie sich wie alle anderen physikalischen Objekte:

\begin{itemize}
    \item \textbf{Verschmelzung:} Wenn \( M_1 + M_2 \leq M_{\text{BEC}} \), verschmelzen sie zu einem neuen BEC-Objekt.
    \item \textbf{Abprallen:} Wenn die Relativgeschwindigkeit hoch genug ist, prallen sie elastisch auseinander.
    \item \textbf{Zerrüttung:} Bei geringer Relativgeschwindigkeit und \( M_1 + M_2 > M_{\text{BEC}} \) zerreißen die Kerne; die Masse wird in die Umgebung überführt.
    \item \textbf{Doppelkern-Systeme:} Die Kerne bleiben getrennt und umkreisen sich in einer gemeinsamen Umgebung.
\end{itemize}

\section{Gravitationswellen von BEC-Objekten}

Verschmelzungen von BEC-Objekten erzeugen Gravitationswellen – ähnlich wie Neutronensterne oder Schwarze Löcher in der ART. Unterschiede:

\begin{itemize}
    \item BEC-Objekte haben keinen Ereignishorizont, sondern eine Oberfläche
    \item Die Wellenform weicht ab, weil die Materie nicht „verschluckt“ wird
    \item Nach der Verschmelzung gibt es Nachglühen
\end{itemize}

Die LIGO/Virgo-Daten sind mit beiden Modellen konsistent, da die beobachteten Massen (\( \sim 10 - 100 M_\odot \)) weit unter \( M_{\text{BEC}} \) liegen.

\section{Zusammenfassung}

Die WDBT+ liefert ein konsistentes, singularitätenfreies Bild der kompakten Objekte:

\begin{itemize}
    \item Die kompaktesten Objekte im Universum sind BEC-Objekte – Bose-Einstein-Kondensate aus Cooper-Paaren, stabilisiert durch das Quantenpotential.
    \item Ihre minimale Größe ist die fundamentale Längenskala \( l_0 \approx 1,8 \times 10^{-15} \) m.
    \item Ihre maximale Masse ist \( M_{\text{BEC}} \approx 3 \times 10^6 M_\odot \).
    \item Beobachtete supermassereiche Objekte (bis \( 10^{10} M_\odot \)) sind Systeme aus einem zentralen BEC-Kern und einer massiven Umgebung.
    \item Bei Kollisionen verhalten sich BEC-Objekte wie normale Objekte: Verschmelzen, Abprallen, Zerrütten.
\end{itemize}